% SEGMENT OLP-0677-S01
% Part: proof-theory
% Chapter: normalization
% Section: segments

\documentclass[../../../include/open-logic-section]{subfiles}

\begin{document}

\olfileid{pt}{nor}{seg}

\olsection{தொடர்களும் வெட்டுகளும்}

!!^a{introduction} விதியைத் தொடர்ந்து !!a{elimination} விதி
வருவது !!a{proof} இல் ஒரு சுற்றுவழி; !!a{proof} ஐ
இயல்பாக்குவது இத்தகைய சுற்றுவழிகளை அகற்றுவதாகும்.
ஆனால் \Elim{\lor}, \Elim{\lexists} விதிகளால்,
“நாம் அகற்ற விரும்பும் சுற்றுவழி” என்பதையே
வரையறுப்பது சற்றுச் சிக்கலாகிறது. சிக்கல் இதுதான்:
இந்த விதிகள் இருப்பதால், ஒரு சுற்றுவழியில்
!!{introduction} விதியை !!{elimination} விதி
\emph{உடனடியாகத்} தொடரும் என்று உறுதியில்லை.
இடையில் \Elim{\lor} அல்லது \Elim{\lexists}
வரலாம்; எடுத்துக்காட்டாக,
\[
\AxiomC{}
\DeduceC{$\lexists[x][!C(x)]$}
\AxiomC{$\Discharge{!A}{x}$}
\DeduceC{$!B$}
\DischargeRule{\Intro{\lif}}{x}
\UnaryInfC{$!A \lif !B$}
\DischargeRule{\Elim{\lexists}}{y}
\BinaryInfC{$!A \lif !B$}
\AxiomC{}
\DeduceC{$!A$}
\RightLabel{\Elim{\lif}}
\BinaryInfC{$!B$}
\DisplayProof
\]

% SEGMENT OLP-0677-S02
உண்மையில், \Intro{\lif} ஐயும் \Elim{\lif} ஐயும்
எத்தனை \Elim{\lor} அல்லது \Elim{\lexists}
விதிகள் வேண்டுமானாலும் பிரிக்கலாம்; எடுத்துக்காட்டாக,
\[
\AxiomC{}
\DeduceC{$\lexists[x][!C(x)]$}
\AxiomC{}
\DeduceC{$!D \lor !E$}
\AxiomC{$\Discharge{!A}{x}$}
\DeduceC{$!B$}
\DischargeRule{\Intro{\lif}}{x}
\UnaryInfC{$!A \lif !B$}
\AxiomC{$\Discharge{!A}{x}$}
\DeduceC{$!B$}
\DischargeRule{\Intro{\lif}}{x}
\UnaryInfC{$!A \lif !B$}
\RightLabel{\Elim{\lor}}
\TrinaryInfC{$!A \lif !B$}
\DischargeRule{\Elim{\lexists}}{y}
\BinaryInfC{$!A \lif !B$}
\AxiomC{}
\DeduceC{$!A$}
\RightLabel{\Elim{\lif}}
\BinaryInfC{$!B$}
\DisplayProof
\]
ஒன்றுக்கு மேற்பட்ட~\Intro{\lif} களைத் தொடர்ந்து
வருவதால், ஒரே \Elim{\lif} கூட பல சுற்றுவழிகளில்
இடம்பெறலாம் என்பதை இந்த எடுத்துக்காட்டு காட்டுகிறது.
இந்த எல்லா வாய்ப்புகளையும் கணக்கில் கொள்ள,
\Log{N1i} நிறுவலில் !!{formula} நிகழ்வுகளின்
சில வகை வரிசைகளை \emph{தொடர்கள்} என்று கொண்டு,
அவற்றின் வழியாகச் சுற்றுவழிகளை வரையறுக்கிறோம்.
எடுத்துக்காட்டில் அது~$!A \lif !B$ இன் நிகழ்வுகளால்
ஆன வரிசை: $\Elim{\lor}$ அல்லது
$\Elim{\lexists}$ இன் சிறுமுன்கோளாக இருக்கும்
ஒரு நிகழ்வை அந்த அனுமானத்தின் முடிவாக இருக்கும்
அடுத்த நிகழ்வு தொடர்கிறது. இவ்வெடுத்துக்காட்டில்
அத்தகைய இரண்டு தொடர்கள் உள்ளன; ஒவ்வொன்றிலும்
மூன்று~$!A \lif !B$ நிகழ்வுகள் உள்ளன.
இப்போது முறையான வரையறை:

% SEGMENT OLP-0677-S03
\begin{defn}
\Log{N1i} !!{proof}~$\delta$ இல் \emph{தொடர்}
என்பது, ஒரே !!{formula}~$!A$ இன்
$!A_1$, \dots, $!A_n$ நிகழ்வுகள்~$\delta$ இல்
அமைக்கும் வரிசை;
அதற்கு பின்வரும் நிபந்தனைகள் உண்டு:
\begin{enumerate}
  \item $!A_1$ என்பது \Elim{\lor} அல்லது
  \Elim{\lexists} இன் முடிவாக இருக்கக்கூடாது;
  \item $!A_n$ என்பது \Elim{\lor} அல்லது
  \Elim{\lexists} இன் சிறுமுன்கோளாக
  இருக்கக்கூடாது;
  \item $i = 1$, \dots,~$n-1$ ஒவ்வொன்றுக்கும்,
  $!A_i$ என்பது \Elim{\lor} அல்லது
  \Elim{\lexists} இன் சிறுமுன்கோள்;
  $!A_{i+1}$ என்பது அதே அனுமானத்தின் முடிவு.
\end{enumerate}
அந்தத் தொடரின் \emph{நீளம்}~$n$.
\end{defn}

% SEGMENT OLP-0677-S04
ஒவ்வொரு தொடரும் அகற்றப்பட வேண்டிய சுற்றுவழி
அல்ல. அகற்றப்பட வேண்டியவை \emph{வெட்டுகள்}
எனப்படும்.

\begin{defn}
\emph{வெட்டு} என்பது $!A_n$ ஓர்
!!a{elimination} அனுமானத்தின் முதன்முன்கோளாக
இருக்கும் ஒரு தொடர்; மேலும், $n=1$ ஆகவும்
$!A_1 = !A_n$ ஓர் !!a{introduction}
அனுமானத்தின் அல்லது~\FalseInt இன் முடிவாகவும்
இருக்க வேண்டும், அல்லது~$n>1$ ஆக இருக்க வேண்டும்.
\end{defn}

வெட்டான தொடர் ஓர் !!a{introduction} விதியின்
முடிவில் தொடங்க வேண்டியதில்லை. அது ஓர்
!!a{elimination} விதியின் முதன்முன்கோளில்
முடிவதே பொதுத் தேவை. ஆனால் நீளம்~$1$
கொண்ட தொடர் வெட்டாக இருக்க வேண்டுமானால்,
அது !!{introduction} விதியின் முடிவாகவோ
பொய்யிலிருந்து வருவித்தலின் முடிவாகவோ இருக்க வேண்டும்.

% SEGMENT OLP-0677-S05
\begin{defn}
ஒரு வெட்டின் \emph{நிலை}~$\depth{!A}$.
!!a{proof}~$\delta$ இன் \emph{வெட்டு
நிலை}~$\cutrank{\delta}$ என்பது~$\delta$
இல் உள்ள வெட்டுகளின் நிலைகளில் மிகப் பெரியது.
ஒரு வெட்டின் நிலை~$\cutrank{\delta}$
ஆக இருந்தால், அதாவது அதன் நிலை மிகப் பெரியதாக
இருந்தால், அது~$\delta$ இல்
\emph{மிகுநிலை} வெட்டு. $\delta$ இன்
\emph{வெட்டு நீளம்} என்பது அதன் எல்லா
மிகுநிலை வெட்டுகளின் நீளங்களின் கூட்டுத்தொகை.
\end{defn}

\end{document}
